%
% einleitung.tex -- Beispiel-File für die Einleitung
%
% (c) 2020 Prof Dr Andreas Müller, Hochschule Rapperswil
%
% !TEX root = ../../buch.tex
% !TEX encoding = UTF-8
%
\section{Zusammenfassung}
\kopfrechts{Zusammenfassung}

Diese Arbeit untersucht, inwiefern sich die klassische Faktorisierung von Polynomen
\index{Faktorisierung}%
über ihre Nullstellen auf analytische Funktionen verallgemeinern lässt.
Ausgehend von Eulers heuristischer Produktdarstellung des Sinus,
\index{Produktdarstellung}%
die als Nebenprodukt eine elegante Lösung des Basel-Problems liefert,
\index{Basel-Problem}%
wird der Weg zu Weierstraß' rigorosem Produktsatz für ganze Funktionen entwickelt.
\index{Weierstrass@Weierstraß}
Ein eigener Abschnitt widmet sich der Gamma-Funktion, ihrer Produktdarstellung
und der berühmten Reflexionsformel, die Sinus- und Gammafunktion auf elegante Weise verknüpft.
Abschließend werden praktische Anwendungen in Signalverarbeitung, Zahlentheorie
\index{Signalverarbeitung}%
\index{Zahlentheorie}%
und Physik skizziert.
\index{Physik}%

\section{Motivation und Fragestellung}
\kopfrechts{Motivation und Fragestellung}

Ein Polynom $p(x)$ vom Grad $n$ ist bis auf einen skalaren Vorfaktor vollständig
durch seine $n$ Nullstellen $x_1, \ldots, x_n$ bestimmt.
\index{Nullstellen}%
Die Faktorisierung
\begin{equation}
p(x) = a_n (x - x_1)(x - x_2) \cdots (x - x_n)
\end{equation}
ist ein Grundpfeiler der Algebra.
Die naheliegende Frage lautet: Gilt Ähnliches auch für transzendente Funktionen
wie den Sinus?
Die Sinusfunktion besitzt unendlich viele Nullstellen, nämlich alle ganzzahligen
Vielfachen von $\pi$.
Kann man $\sin(x)$ als unendliches Produkt über die zugehörigen Linearfaktoren schreiben?

Diese Frage führt direkt zu einer der schönsten Formeln der Analysis,
die Leonhard Euler 1734 aufstellte, und schließlich zum allgemeinen
\index{Euler, Leonhard}%
Weierstraß-Produktsatz der Funktionentheorie.
Zugleich liefert sie einen überraschend direkten Beweis des Basel-Problems.
Darüber hinaus zeigt sich, dass dieselben Ideen die Gamma-Funktion,
\index{Gamma-Funktion}%
die natürliche Verallgemeinerung der Fakultät, zutiefst mit der Theorie
\index{Fakultät}%
der ganzen Funktionen verknüpfen.

